% ---------------------------------------------------------------------------
% Replacement Section 6 for "Adaptive SIMEX for Poisson Regression with
% Classical Measurement Error" --- NHANES 2017-2018 application.
% Requires: nhanes_table.tex, fig4_sensitivity.*
% Reproducible via code/nhanes_analysis.py (data downloaded from CDC).
% ---------------------------------------------------------------------------

\section{Real Data Application: Self-Reported Weight and Prescription
Medication Use}\label{sec:app}

The application of this section is one in which the measurement error
is directly measurable: the relation between body weight and
prescription medication use in the National Health and Nutrition
Examination Survey (NHANES) 2017--2018.
NHANES is its own validation study: every examined participant reports
weight in a household interview (\texttt{WHD020}) \emph{and} is weighed
in the mobile examination center (\texttt{BMXWT}), so the classical
error variance can be estimated from the same survey rather than
asserted. Count outcomes with a plausibly weight-linked intensity are
scarce in survey data; the number of prescription medications recorded
in the prescription medication file (\texttt{RXDCOUNT},
\citep{nhanes-rxq}) is a genuine count and is the outcome
$Y$ throughout.

After dropping records with missing or implausible values (self-reported
or measured weight outside $(30,350)$\,kg, missing covariates), the
analysis sample is $n=5{,}628$ complete cases. The count model is
\begin{equation}
  \log\E[Y_i\mid X_i] = \beta_0 + \beta_w W_i/10
    + \beta_1\,\mathrm{age}_i/10 + \beta_2\,\mathrm{female}_i
    + \beta_3\,\mathrm{height}_i/100\mathrm{m}
    + \beta_4\,\mathrm{income}_i,
  \label{eq:nhanes}
\end{equation}
where $W_i$ is self-reported weight in kg (the error-prone covariate;
all other covariates (age, sex, measured height, and the family
income-to-poverty ratio) are treated as error-free). The self-report
error $W_i-\mathrm{BMXWT}_i$ has mean $-0.71$\,kg (under-reporting) and
variance $\hat\sigma^2_\epsilon = 24.1$\,kg$^2$ (s.d.\ $4.9$\,kg),
which we feed to all correction methods; because the mean difference is
small relative to the error variance the classical zero-mean
approximation is mild, and Section~\ref{sec:app-sens} varies
$\sigma^2_\epsilon$ over a fourfold range. Weight enters per 10\,kg.
For comparison, the \emph{measured-weight fit}, which replaces
self-reported by measured weight (observed for every participant), is shown
as the gold-standard anchor, a check unavailable in typical SIMEX
applications.

Table~\ref{tab:nhanes} reports the weight coefficients. The naive fit
gives $+1.189$ per 10\,kg (s.e.\ $0.076$). The attenuation implied by
$\hat\sigma^2_\epsilon$ relative to the population weight variance
(about $5\%$) is modest, and accordingly all correction methods move
the estimate only slightly: regression calibration $+1.250$, corrected
score $+1.253$, standard SIMEX $+1.253$, Adaptive SIMEX (1-SE rule)
$+1.247$, and Adaptive SIMEX (argmin) $+1.214$, all with s.e.\
$\approx0.081$. The measured-weight anchor is $+1.095$ (s.e.\ $0.072$),
about $2$ standard errors \emph{below} the corrected estimates. Two
readings of this gap are consistent with our theory. First, the
correction estimates the coefficient of a \emph{noise-free} latent
weight on the self-report scale, whereas the anchor uses the measured
weight, a different (equally noisy but differently distributed)
regressor; the classical model asserts their coefficients coincide, and
a shortfall is the signature of nonclassical error: heavier
participants under-report by more, inducing correlation between
$\epsilon$ and $X$ that all classical-error methods, including ours,
will not remove. Second, the agreement among the four corrections
(similar to the simulations of Section~\ref{sec:sims1} at small error)
indicates that the correction machinery itself is stable at
$n=5{,}628$; the CV layer selects minimal extrapolation
($(\lambda_{\max},\text{degree})=(0.5,1)$), consistent with
Section~\ref{sec:sims1}'s finding that held-out deviance tracks the
naive pseudo-parameter.

\begin{table}[t]
\centering
\small
\caption{NHANES 2017--2018 application: Poisson regression of the
number of prescription medications on body weight (coefficient per
10\,kg, robust standard errors). The error-prone covariate is
self-reported weight; the classical error variance is estimated
from the self-report minus measured differences
($\hat\sigma^2_\epsilon = 24.1\,\mathrm{kg}^2$,
mean difference $-0.7\,\mathrm{kg}$). The
measured-weight fit is shown as the gold-standard anchor.}
\label{tab:nhanes}
\begin{tabular}{lccc}
\toprule
Method & $\hat\beta$ (10\,kg) & SE & Selected $(\lambda_{\max},$ degree$)$ \\
\midrule
Naive (self-reported weight) & +1.189 & (0.076) & -- \\
Regression calibration & +1.250 & (0.080) & -- \\
Corrected score (Nakamura) & +1.253 & (0.081) & -- \\
Standard SIMEX & +1.253 & (0.081) & -- \\
Adaptive SIMEX (1-SE rule) & +1.247 & (0.081) & ((, n) \\
Adaptive SIMEX (argmin) & +1.214 & (0.078) & ((, n) \\
Measured weight (anchor) & +1.095 & (0.072) & -- \\
\bottomrule
\end{tabular}
\end{table}

\begin{figure}[t]
\centering
\includegraphics[width=.8\textwidth]{fig7_nhanes_forest.pdf}
\caption{NHANES 2017--2018: weight coefficients (per 10\,kg) with
95\% confidence intervals by method. Dashed line: measured-weight
anchor.}\label{fig:forest}
\end{figure}

Figure~\ref{fig:forest} displays the estimates graphically.

\subsection{Sensitivity analysis over the error variance}
\label{sec:app-sens}

Because $\hat\sigma^2_\epsilon$ is itself estimated (and the classical
assumption is at best approximate), we refit every correction method
with $\sigma^2_\epsilon = m\,\hat\sigma^2_\epsilon$ for
$m\in\{0.25,0.5,\dots,2\}$, following the original study plan. Figure~\ref{fig:sens} plots the weight
coefficients: the standard SIMEX, Adaptive SIMEX, and corrected-score
estimates rise from $+1.204$ ($m=0.25$) to $+1.305$--$+1.323$
($m=2$), a range of about one and a half standard errors across a fourfold
misstatement of the error variance, bracketing the measured-weight
anchor. The application therefore illustrates both the use and the
honest limitation of the method: with directly measurable error the
corrections are stable and mutually consistent, and the residual
anchor gap points to the error structure rather than the estimator.

\begin{figure}[t]
\centering
\includegraphics[width=.62\textwidth]{fig4_sensitivity.pdf}
\caption{NHANES application: weight coefficient (per 10\,kg) as a
function of the error-variance multiplier
$m$ ($\sigma^2_\epsilon = m\hat\sigma^2_\epsilon$). Dashed: measured
weight (anchor); dotted: naive.}\label{fig:sens}
\end{figure}

% Notes to the author (remove before compiling):
%  - \citep{nhanes-rxq}: add the NHANES 2017-2018 RXQ_RX documentation
%    citation (National Center for Health Statistics) to the references.
%  - \ref{sec:sims1} refers to Section 5.1 of the replacement Section 5.
% ---------------------------------------------------------------------------
